\documentclass[11pt]{article}
\usepackage[T1]{fontenc}
\usepackage{lmodern,amsmath,amssymb,amsthm,mathtools}
\usepackage[margin=0.86in]{geometry}
\usepackage{booktabs,longtable,array,enumitem,xcolor,fancyhdr,hyperref,xurl}
\hypersetup{colorlinks=true,linkcolor=blue!45!black,urlcolor=blue!45!black,
pdftitle={Prime Squares: Initial Formalization Blueprint and Checked Finite Core},
pdfauthor={Prepared for Hamed Mousavi}}
\definecolor{ink}{HTML}{16324F}
\definecolor{muted}{HTML}{526477}
\pagestyle{fancy}\fancyhf{}
\fancyhead[L]{\small\color{muted}Prime squares: initial formalization blueprint}
\fancyhead[R]{\small 20 October 2026}
\fancyfoot[C]{\small\thepage}
\setlength{\headheight}{14pt}
\setlength{\tabcolsep}{5pt}
\setlength{\parindent}{0pt}\setlength{\parskip}{6pt}
\setlist{nosep,leftmargin=1.5em}
\newcommand{\R}{\mathbb R}\newcommand{\Z}{\mathbb Z}\newcommand{\N}{\mathbb N}
\newcommand{\T}{\mathbb T}\newcommand{\PP}{\mathbb P}
\newcommand{\one}{\mathbf 1}\newcommand{\abs}[1]{\left|#1\right|}
\newcommand{\norm}[1]{\left\|#1\right\|}
\newcommand{\lean}[1]{\nolinkurl{#1}}
\newcommand{\status}[1]{\textbf{Status:} #1.}
\newtheorem{lemma}{Lemma}[section]
\newtheorem{proposition}[lemma]{Proposition}
\newtheorem{definition}[lemma]{Definition}
\newcommand{\page}{\newpage}
\begin{document}
\begin{center}
{\LARGE\bfseries\color{ink}Prime Squares}\\[6pt]
{\Large Initial Formalization Blueprint}\\[4pt]
{\large and a checked finite core}\\[18pt]
For Hamed Mousavi\\
\emph{On Equidistribution of Prime Squares and Its Applications}\\[10pt]
20 October 2026
\end{center}
\vspace{12pt}
\textbf{What has been done.} This package contains a first dependency blueprint for the
attached 96-page manuscript and a small Lean project. The project formalizes finite
exceptional-set inequalities, elementary arithmetic used in the paper, and positivity
of the elementary bracket in the last sieve comparison. These are useful components
of a larger development; none is a formalization of a headline theorem of the paper.

\textbf{What this document does not claim.} It is not an expanded proof of every lemma,
and it is not a certificate for the whole paper. In particular, the circle-method
approximation, high--low estimates, stopping-time proof, Ionescu--Wainger argument,
ergodic divergence argument and sieve estimates remain to be encoded and checked.
The word \emph{target} below always denotes unfinished work.

\textbf{Source lock.} Numbering and page references refer only to the uploaded file
\lean{Prime_Squares (32)(1).pdf}. Its SHA-256 is recorded in
\lean{PROVENANCE.json}. A revised manuscript must be compared with this snapshot
before the paper-to-Lean map is reused.

\textbf{Reproducibility.} Lean is pinned to 4.34.1. Mathlib is pinned to commit
\lean{d13f23b723b8a846827a245b89c10fc7d3f11612}. The source archive includes a lockfile,
the Lean modules, build instructions and scripts for checking proof dependencies.
The verification logs, rather than this prose, record what was actually run.

\textbf{Reading route.} Sections 1--3 define the scope and targets. Sections 4--5 give
complete mathematical explanations of the checked components. Sections 6--11 map
the unfinished arguments. Sections 12--14 record imports, audit points and the next
inputs needed. The appendix inventories numbered statements in the attached PDF.

\page
\tableofcontents

\page
\section{The trust boundary}
There are three different objects in a formalization project.
\begin{enumerate}
\item A \textbf{blueprint} specifies mathematical objects, statements, dependencies and
proof steps. It is written for humans and can itself contain mistakes.
\item A \textbf{Lean proof} is an actual term accepted by Lean's kernel. Elaborating
a statement containing a placeholder is not the same as proving it.
\item A \textbf{conditional formalization} proves a result from explicitly declared
external statements. Its mathematical scope depends on the accuracy and consistency
of those statements as well as on the translation of the final theorem.
\end{enumerate}

For the present starter, no extra analytic-number-theory axioms are declared.
Theorems in the starter have explicit hypotheses for the errors and main terms they
use. Those hypotheses are not assertions that the required analytic estimates have
already been established for prime squares.

For the completed project it is reasonable to import deep published results.
Each import should have one stable declaration, a precise source, the full range of
parameters, and a statement of constant dependence. The project should continue
to prove the new arguments of this paper. Assuming its own high--low estimate or
its final sieve comparison would leave that part of the paper unverified.

An axiom audit must examine the \emph{transitive} dependencies of each headline
theorem, not just count explicit occurrences of \lean{axiom}. Standard Lean/Mathlib
dependencies can include \lean{propext}, \lean{Classical.choice} and
\lean{Quot.sound}. An unfinished proof introduces \lean{sorryAx}. The starter's
scripts reject unapproved dependencies and scan the project sources for proof
placeholders and new axioms.

\subsection*{What the Hua example establishes}
The supplied Hua blueprint contains both a current-paper development and substantial
material from an earlier, broader draft. Its total length is therefore not a useful
measure of how long the prime-squares blueprint should be.

At the repository snapshot inspected for this request, the trusted external
declarations are \lean{Hua.siegel_walfisz} and \lean{Hua.hua_theorem10}. The second
is an exponential-sum estimate, rather than the statement that sufficiently large
admissible integers are sums of five prime squares. A static scan of its default
library closure passed. That scan does not replace an independent rebuild or an
audit of the meanings of its axioms and headline statements. This starter does not
import that repository.

\page
\section{Objects and final statements to encode}
\status{Specifications only; the headline statements are not yet Lean declarations}

Use natural numbers for prime parameters and integer indices for discrete
functions. Define, for \(N\ge2\),
\[
 A_{N,\Lambda}f(x)=\frac1N\sum_{1\le n\le N}\Lambda(n)f(x-n^2),\quad
 A_{N,\square}f(x)=\frac1N\sum_{1\le n\le N}f(x-n^2),
\]
\[
 A_{N,\PP^2}f(x)=\frac1{\pi(N)}\sum_{p\le N}f(x-p^2).
\]
The first and third operators differ in both weights and normalization. In the
divergence theorem the parameter counts the first \(N\) primes instead of bounding
their size. These conventions must not be silently identified.

Put \(L(t)=2+|\log t|\), \(L_2(t)=2+\log(2+|\log t|)\), and
\[
 \Phi_c(t)=\exp\!\left(c\frac{L(t)^{1/2}}{L_2(t)}\right)\quad(t>0).
\]
In bilinear forms define \(t^{2/3}\Phi_c(t)\) to be zero at \(t=0\).
For \(t\ge0\), define
\[
 \Psi_C(t)=\int_0^t
 \exp\!\left(C\sqrt{\log(e+u)\log(e+\log(e+u))}\right)\,du.
\]
The square root in this last expression covers the \emph{product} of the logarithms.

\textbf{Theorem 1.1.} For every \(\varepsilon>0\), there is \(C_\varepsilon\) such
that the displayed restricted improving inequality holds for every allowed \(N,F,G\)
with coefficient \(4\sqrt3/9+\varepsilon\). The \(o(1)\) in the paper is thus a
quantifier over \(\varepsilon\), not an unspecified function added to the theorem.
Strong improving bounds for \(3/2<p\le2\) are separate corollaries.

\textbf{Theorem 1.3.} Encode the endpoint sparse form and then the stated interior
triangle of exponents. Sparsity uses subsets \(E_I\subset I\) of cardinality at
least \(|I|/10\), pairwise disjoint over the family. Cardinalities, dilations and
normalizations must be explicit.

\textbf{Theorem 1.5.} Quantify over every nonatomic ergodic measure-preserving
system and every infinite set of indices \(S\). The conclusion is existence of an
integrable function whose averages along \(N\in S\) fail to converge on a set of
positive measure. Do not weaken this to one system or to the full sequence only.

\textbf{Theorem 1.6.} Define \(\mathcal P\) as the primes represented by
\(a^2+b^2+1\) with \((a,b)=1\). For every \(A>0\), specify constants \(C_A,X_A\)
so that the nonrepresentable \(1\le n<X\), \(n\equiv12\pmod{24}\), have cardinality
at most \(C_AX(\log X)^{-A}\) for \(X\ge X_A\).

\textbf{Section 7.} Theorem 7.2 and Corollaries 7.6 and 7.42 are additional targets:
the modular estimate, quantitative \(L^p\) growth, and almost-everywhere and \(L^1\)
convergence on the stated Orlicz class. They are not substitutes for Theorem 1.5.

\page
\section{The dependency architecture}
\status{Initial map, not a claim that the indicated modules exist already}

\begin{longtable}{p{.15\linewidth}p{.36\linewidth}p{.40\linewidth}}
\toprule
Branch & Principal nodes in the paper & Output and main dependencies\\\midrule
Arithmetic & Section 3, especially 3.1, 3.12, 3.14, 3.20, 3.21, 3.24, 3.27 &
Cohen correlations, kernel bounds, Gauss sums and exceptional-conductor bounds.\\
Improving & 4.1, 4.15, 4.19 &
Theorem 1.1; uses Section 3 and the decomposition in Section 6.\\
Sparse & 5.1, 5.10, 5.17 &
Theorem 1.3; needs the local low estimate and maximal, rather than fixed-scale,
high estimate.\\
High--low & 6.5, 6.9, 6.14, 6.16, 6.19, 6.21 &
Proposition 4.19; uses primes in progressions with an exceptional term, minor-arc
cancellation and multifrequency estimates.\\
Near \(L^1\) & 7.8, 7.15, 7.21, 7.28, 7.2 &
Corollaries 7.6 and 7.42; uses quantitative Ionescu--Wainger/sampling inputs and
the Section 6 model.\\
Divergence & 8.1, 8.6, 8.8 and the final assembly &
Theorem 1.5; uses a character-sum bound, Poisson spacings, fixed-modulus prime
equidistribution and the divergence criterion.\\
Goldbach & 9.2--9.10, Section 10, 11.13, 11.19, 9.17, 9.20, 9.22, 9.24 &
Theorem 1.6; circle-method mean-square approximation, two-stage sieving,
positive main term and exceptional-set assembly.\\\bottomrule
\end{longtable}

The high--low branch is shared by the improving, sparse and near-\(L^1\) results.
The divergence branch shares arithmetic foundations but is logically separate from
the maximal inequality. The Goldbach branch shares exponential sums and local
correlations, but introduces its own mean-square approximation and sieve machinery.

\subsection*{Proposed eventual organization}
Keep declarations under \lean{PrimeSquares}. Suggested future directories are
\lean{Arithmetic}, \lean{HighLow}, \lean{Improving}, \lean{Sparse},
\lean{Maximal}, \lean{Ergodic}, \lean{Goldbach} and \lean{External}.
Maintain one table linking a manuscript label to each Lean declaration. Use stable
LaTeX labels once the sources are available; printed theorem numbers alone will
change during proofreading.

For each node record: the exact statement; normalization and parameter ranges;
dependencies; an expanded proof; Lean name; external assumptions; build status.
The appendix currently gives an inventory, not this complete node-by-node expansion.

\page
\section{A checked finite exceptional-set argument}
\status{Implemented in \lean{PrimeSquares/FiniteAssembly.lean}}

Let \(S\) be a finite set, let \(e_1,e_2:S\to\R\), and let \(t>0\). Define
\[
 B_t=\{n\in S:|e_1(n)|+|e_2(n)|>t\}.
\]
No number-theoretic interpretation is part of these definitions.

\begin{lemma}[Finite energy estimate]
If \(\sum_{n\in S}e_i(n)^2\le B_i\) for \(i=1,2\), then
\[
 |B_t|\le\frac{2(B_1+B_2)}{t^2}.
\]
\end{lemma}
\begin{proof}
For real \(a,b\), the nonnegativity of \((|a|-|b|)^2\) implies
\((|a|+|b|)^2\le2(a^2+b^2)\). On \(B_t\), therefore,
\(t^2\le2(e_1(n)^2+e_2(n)^2)\). Sum this inequality on \(B_t\), enlarge
the nonnegative square sum to \(S\), and divide by \(t^2>0\).
\end{proof}

\begin{lemma}[Positivity outside the exceptional set]
Let \(R(n)\) be a nonnegative integer and suppose that
\[
 R(n)\ge M(n)+e_1(n)-e_2(n)-u(n),\quad
 M(n)\ge4t,\quad u(n)\le t.
\]
If \(n\notin B_t\), then \(R(n)\ge2t>0\).
\end{lemma}
\begin{proof}
Use \(e_1-e_2\ge-|e_1|-|e_2|\) and \(|e_1|+|e_2|\le t\). The lower
bound becomes \(4t-t-t=2t\).
\end{proof}

Combining the two lemmas gives the actual Lean result
\lean{PrimeSquares.zero_count_card_bound}:
\[
 \abs{\{n\in S:R(n)=0\}}\le\frac{2(B_1+B_2)}{t^2}.
\]
The hypotheses about \(M,R,u,e_i\) are part of its statement. In the paper's final
assembly take \(M=\Delta X_n\Theta_n(z)/2\),
\(t=cN(\log N)^{-C_{A_0}}\), and \(e_i=E_i\).
If each mean square is at most \(C_iN^4(\log N)^{-K}\), the bound becomes
\[
 \abs{\{R=0\}}\le
 \frac{2(C_1+C_2)}{c^2}N^2(\log N)^{-K+2C_{A_0}}.
\]
Choosing \(K\) large enough yields the desired logarithmic saving for this finite
range. The analytic estimates, eventual parameter choices, identification of \(R\)
with actual representations, and treatment of the initial short range remain
formalization work. Thus this is the finite last step of the Goldbach proof, not
Theorem 1.6 itself.

\page
\section{Other checked components}
\status{Implemented in \lean{LocalArithmetic.lean} and \lean{SieveConstant.lean}}

\subsection*{Congruences and the shared-prime observation}
For every natural number \(n\) with \(n\equiv1\pmod2\) and \(n\not\equiv0\pmod3\),
the module proves \(n^2\equiv1\pmod{24}\). The proof reduces to the 24 possible
residues; Lean checks the finite case split. For a prime \(p>3\), primality excludes
divisibility by 2 and 3, so the same conclusion follows.

Consequently, if \(r\equiv11\pmod{24}\) and \(p>3\) is prime, then
\(r+p^2\equiv12\pmod{24}\). The restriction \(p>3\) is included explicitly.
The starter also proves the shift equivalences
\[
 m=a^2+b^2+p^2\iff m+1=a^2+b^2+p^2+1,
 \qquad m\equiv11\pmod{24}\iff m+1\equiv12\pmod{24}.
\]
These clarify the relation between the abstract's formulation and Theorem 1.6;
they do not establish an exceptional-set estimate.

Suppose \(p,m\) are prime, \(n=m^2+r\), and \(p\mid n,r\). Then \(p\mid m^2\),
so \(p\mid m\); primality of \(m\) forces \(p=m\). This is the finite divisibility
observation used in Lemma 11.2. The divisor-counting estimate and the full sieve
comparison in that lemma are not yet formalized.

\subsection*{The elementary bracket in Lemma 9.24}
The manuscript reduces its final sieve comparison to
\[
 \Delta\ge\frac{e^{\gamma/2}}{\sqrt\pi}
 \left(\frac{2\sqrt2}{3}\log2-\frac{27}{64}\right).
\]
Lean proves positivity of the parenthesized expression. There is a particularly
short rigorous proof which does not need decimal approximations:
\[
 \sqrt2\ge\frac43,\qquad \log2\ge1-\frac12=\frac12,
\]
where the second inequality is the standard \(1-x^{-1}\le\log x\) for \(x>0\).
Both lower bounds are positive, and therefore
\[
 \frac{2\sqrt2}{3}\log2-\frac{27}{64}
 \ge\frac49-\frac{27}{64}=\frac{13}{576}>0.
\]
The declaration \lean{sieve_bracket_positive} proves this conclusion. The separate
declaration \lean{scaled_sieve_bracket_positive} allows multiplication by any
strictly positive real factor.

This does \emph{not} verify the integral evaluation for \(f_{1/2}(9/8)\), the upper
bound on \(\mathcal W(9/8)\), the factor \(e^{\gamma/2}/\sqrt\pi\), or the stronger
numerical inequality \(\Delta>0.17\). Those must be added before claiming Lemma 9.24.

\page
\section{Improving estimates: the scalar optimization and its inputs}
\status{Unfinished target; the following is the blueprint for Section 4}

Separate the prime parameter from the spatial scale. For prime squares, a sum up
to \(N\) has support at scale \(N^2\); the abstract theorem uses a spatial-scale
variable. A typed name such as \lean{primeScale} versus \lean{spaceScale} prevents
an incorrect substitution of \(N\) for \(N^2\).

Theorem 4.1 first supplies a comparison operator \(B_N\). Applied to indicators,
its decomposition gives, with \(\delta=|F||G|/N^2\),
\[
 N^{-1}|\langle B_N\one_F,\one_G\rangle|
 \lesssim\delta J^s(\log J)^\gamma+
 \delta^{1/2}J^{-r}(\log J)^\beta.
\]
The optimal \(J\), up to bounded factors, is
\[
 J=\delta^{-1/(2(s+r))}(\operatorname{Log}\delta)^{(\beta-\gamma)/(s+r)}.
\]
Prove first the estimate when this \(J\) is in the permitted range. If it is too
large, use the elementary local kernel estimate. This second branch is essential:
choosing an inadmissible \(J\) is not an application of the decomposition.

For the prime operator the corresponding bound is
\[
 \delta Q^s(\log Q)^\alpha+
 \delta^{1/2}Q^{-r}
 \exp\!\left(c_1\frac{(\log Q)^\kappa}{\log\log Q}\right)
 +\text{error}.
\]
When \(\operatorname{Log}\delta\) is small enough relative to \((\log N)^c\),
take the dyadic \(Q\) specified by (4.14). Prove separately that it is admissible,
that rounding changes the estimate by a bounded factor, and that the error is
absorbed. In the complementary range, transfer the comparison estimate using the
pointwise domination by \(B_N\). This yields
\[
 \delta^{1-s/(2(s+r))}\Phi_{c_*,\kappa}(\delta),\qquad
 c_* = \frac{c_1s}{2^\kappa(s+r)^{\kappa+1}}+o(1).
\]
For \(s=1,r=1/2,\kappa=1/2,c_1=2+\varepsilon\), the critical power is \(2/3\)
and the limiting coefficient is \(4\sqrt3/9\).

\textbf{Separate Lean tasks.} Establish the finite bilinear norm inequalities;
formalize the elementary optimization including the two parameter ranges; prove
the dyadic-to-full annular summation; and then prove restricted-to-strong
interpolation for \(p>3/2\). None follows just from writing the numerical value
of the coefficient into a declaration.

\page
\section{Sparse bounds: local control and recursion}
\status{Unfinished target; Sections 5 and the stopping-time conventions in Section 2}

The sparse proof needs more than Theorem 1.1. A stopping time \(\tau(x)\) depends
on the spatial point. Its high term must be controlled by a maximal \(\ell^2\)
estimate over all permitted scales. The low estimate must be local enough that
admissibility replaces its local density by the parent density.

The key tasks are these.
\begin{enumerate}
\item Define integer intervals, their dilations, a dyadic grid, and the admissible
stopping times. Record the exact support radius of each kernel.
\item Prove Lemma 5.1 from the local kernel estimate, the local low estimate, the
maximal high estimate and the square-summable errors. Include the small-scale
branch, where the abstract high--low split is not available.
\item Prove Lemma 5.10. Select maximal bad dyadic intervals by the two density
tests. Bound the sum of child lengths by \(|I_0|/4\). Construct disjoint major
subsets and iterate the local inequality.
\item In Lemma 5.17 split according to \(\log\tau\le C_0L^2\) and its complement,
where \(L=\operatorname{Log}(\delta_0^{-1})\). The first range uses the square
averages; the second uses the prime-square decomposition with an admissible
denominator cutoff.
\item Convert the local sparse form with \(cI\) on the input side into the stated
form. Prove the required grid/dilation reduction rather than identifying the
two families by notation.
\end{enumerate}

For finite input sets, begin with a finite set of scales and a finite recursion.
If using an infinite family later, specify its sum and prove the limiting step.
This keeps the initial recursive construction elementary and exposes precisely
which part needs an infinite argument.

There is a useful formalization trap here. An arbitrary real-valued supremum in
Lean uses conditional completeness. Without a boundedness hypothesis, it must
not be treated as a maximum with all the usual inequalities. One can first work
with finite truncated maxima, or use extended nonnegative reals for a maximal
function and prove finiteness afterward.

Finally, the strong sparse estimates for exponents in the interior of the triangle
need their own interpolation/convexity argument. They should be separate
declarations depending on the endpoint local bound and the trivial bounds, not
an additional assumed theorem.

\page
\section{The high--low model and exceptional conductors}
\status{Unfinished target; the main shared arithmetic and analytic branch}

At prime scale \(N\), the annular multiplier is approximated by the model \(V_N\)
in (6.8). Its ordinary coefficients are normalized unit quadratic Gauss sums.
If an exceptional character is present, its contribution has the minus sign
displayed in that formula, an induced character modulo each multiple of its
conductor, and a continuous model with parameter \(\beta\).

Use explicit scale-indexed data. For example, define an optional datum at \(N\)
containing \((q_\chi,\chi,\beta)\), and define
\(\mathcal M(\chi)\) to be the scales where that character is selected. Do not
pick one character and use it at all scales. Page uniqueness must be strong
enough to justify the later regrouping by conductor.

\textbf{Arithmetic layer.} Prove the Cohen correlation by CRT, the odd prime-power
evaluation and the explicit 2-adic cases. Build the restricted correlations and
Euler products of Section 3 from these identities. Prove the ordinary Gauss bound
with its subexponential factor, the real-character bound, and the stronger
\(\sigma_\chi B_\chi^{1/2}\ll q_\chi^{-1/25}\) estimate. The last estimate is used
to sum over changing exceptional characters.

\textbf{Approximation layer.} Encode Lemma 6.5 and then Lemma 6.9. The primes-in-
progressions input must supply a single exceptional datum valid over the interval
of partial summation \([N/2,N]\), with uniform errors. A fixed-modulus PNT alone
does not provide this growing-modulus input.

\textbf{Exact decomposition.} The ordinary split is by \(q\le Q\). The exceptional
split is by \(d=q/q_\chi\le Q/B_\chi\), with its own cutoff. Prove the multiplier
identity (6.13) before norm estimates. Empty exceptional sums, including
\(Q/B_\chi<1\), must be covered by the definitions.

\textbf{Low and high bounds.} Proposition 6.14 uses inverse transforms and the
arithmetic kernel identities. Proposition 6.21 needs Lemma 6.16, the denominator-
prefix extension 6.19, the cutoff-difference square functions, and a final sum
over \(\chi\). At each scale there is at most one exceptional datum; the square
sum over characters and Page uniqueness per dyadic conductor range are separate
steps. Only after both bounds and the error estimate are checked does one obtain
Proposition 4.19.

\page
\section{The Orlicz maximal estimate and convergence}
\status{Unfinished target; Section 7}

The quantitative Ionescu--Wainger input must retain its constants. Proposition
7.8 specifies a rational-frequency set containing every denominator at most \(Q\),
a period \(P\), a denominator bound \(B\), a sufficiently small width, and norm
comparisons for a range of exponents. It invokes Tao's Lemma 5.1, Theorems 1.6
and 1.7, and Proposition 1.5. These should have explicit external interfaces if
they are not formalized here.

Next prove the continuous multiplier bounds uniformly in \(0<\beta\le1\)
(Lemma 7.15), followed by the finite-group measure and sampling argument
(Proposition 7.21). The measure is allowed to be complex but has total variation
at most one. Dependence of constants on the period or on \(\beta\) would break
the later optimization.

Proposition 7.28 combines this low model with an \(\ell^2\) high estimate. The
exceptional characters of conductor above \(Q\) vary with scale and are retained
in the high term. The low exceptional character, when it exists, is controlled
using Page uniqueness. This assembly must use the same model and normalization
as Section 6.

For the modular extension, do not stop at the restricted estimate on \(\one_F\).
Let \(C_\infty\) bound the masses of the positive kernels. Remove
\(|f|\le\lambda/(2C_\infty)\). Decompose the remaining function into level sets
\[
 E_j=\{2^j\lambda/(2C_\infty)<|f|\le2^{j+1}\lambda/(2C_\infty)\}.
\]
Allocate thresholds using a positive summable sequence such as
\(a_j=6/(\pi^2(j+1)^2)\). Prove the union inclusion, apply the restricted
estimate to each \(E_j\), and prove that the polynomial factor in \(j\) inside
the threshold can be absorbed by enlarging the coefficient of \(\Psi_C\).
This is a quantitative estimate for this particular Young function, not just
sublinearity. It gives the truncated modular estimate (7.3), then (7.4).

Corollary 7.6 follows by layer cake and a one-variable integral estimate, with
\(\varepsilon=p-1\). Corollary 7.42 first transfers finite truncated discrete
estimates to an invertible probability-preserving system. Prove bounded-function
convergence using the cited polynomial-prime result and control the higher prime
powers there. Abel summation handles the unweighted prime averages. Finally use
Orlicz truncation, doubling and the modular maximal estimate to pass to unbounded
functions; use uniform \(L^1\) bounds for \(L^1\) convergence.

\page
\section{Persistent universal divergence}
\status{Unfinished target; Section 8}

For each modulus \(q\), define
\(\Lambda_q=\{b^2\bmod q:(b,q)=1\}\). The initial finite tasks are CRT
compatibility, inclusion in the prime-square residue set, and, for odd squarefree
\(q\),
\[
 |\Lambda_q|=\frac{\varphi(q)}{2^{\omega(q)}},\qquad
 s_q=\frac q{|\Lambda_q|}=\prod_{p\mid q}\frac{2p}{p-1}.
\]
The inclusion uses Dirichlet's theorem for a unit square root; it must not claim
that every prime residue, including nonunits, is part of \(\Lambda_q\).

Lemma 8.8 expands each quadratic-residue indicator as
\(\tfrac12(1+\chi_p-\one_{\{0\}})\). Terms containing a zero indicator are bounded
by a constant depending on the number of shifts. For the others, apply the Weil
bound to the squarefree product of the distinct shifted linear factors. This
proves the correlation hypothesis for the imported Poisson-spacings theorem.

Choose pairwise disjoint blocks of odd primes with block sizes tending to
infinity, and a disjoint prime sequence \(u_j\). The block products \(q_j\)
satisfy \(\PP_{q_j}(\Lambda_{q_j})\le2^{-j}\), whereas
\(\PP_{u_j}(\Lambda_{u_j})\ge1/3\). The construction must ensure the combined
collection is pairwise coprime.

Theorem 8.6 supplies Poisson spacings. The thickening identity in the proof
expresses the covered proportion as
\[
 \PP_{q_j}(\Lambda_{q_j}^{\gamma})
 =\int_0^\gamma F_j(\theta)\,d\theta+O(s_{q_j}^{-1}).
\]
Dominated convergence gives \(1-e^{-\gamma}\), so the deficit is \(o(\gamma)\).
The pair-separation condition uses LaVictoire's Lemma 3.2 in addition to this
one-gap distribution; it is not an automatic restatement of Poisson convergence.

Finally, fixed-modulus prime equidistribution gives limiting mass
\(1/|\Lambda_Q|\) at every \(a\in\Lambda_Q\). The strict lower bound required
by (8.5) then holds with, for example, \(\beta=1/2\). Apply the full imported
divergence criterion, including its transference/Banach-principle conclusions,
to obtain exactly the persistence and system quantifiers of Theorem 1.5.

If the chosen formal ergodic infrastructure assumes invertibility, a separate
extension/transfer argument is needed to match the unqualified system statement
in Theorem 1.5. Simply adding invertibility would narrow the theorem.

\page
\section{The Goldbach branch}
\status{Only the finite final assembly and small supporting pieces are checked}

The reductions in Section 9 produce a lower sieve count minus an upper correction
count. Their errors must have mean-square bounds with an arbitrarily large
logarithmic saving. The common main term is
\(X_n\Theta_n(z)\); its positivity and uniform lower bound are separate lemmas.

\textbf{Circle method.} Lemma 9.10 is the principal target. Its coefficients are
independent of \(n\), and its conclusion is a mean square of a sum over auxiliary
moduli and \(\ell\). Section 10 supplies the minor-arc bound, the major-arc local
factor \(G_R\), the truncated continuous factor \(H_{\ell,R}\), the tail estimate
and the major-arc error. The same small-prime parameter \(R\) must be retained
in both main-term factors until its tails have been shown to be summable.

\textbf{Sieve.} Lemma 11.2 replaces variable sets of sieve primes by fixed sets,
which is needed before applying Lemma 9.10 to coefficients independent of \(n\).
The exact local densities are computed in Lemmas 11.6 and 11.8. For primes at
most \(R\), the uniform product dimensions are \(2/3\) and \(3/2\); one cannot
apply the desired half-dimensional and linear sieves over the whole range.

Lemma 11.13 removes \(p\le R\) first at level
\(D_0=\exp((\log\log N)^3)\). The remaining level is \(D/D_0\).
Because outer weights have both signs, the choice of preliminary upper or lower
weights must depend on that sign. Prove this finite algebra before estimating
the two stages. The coefficient bounds and independence of \(n\) must survive
the combination. The outer dimensions are then \(1/2\) and \(1\).

Lemma 11.19 performs the weighted sum over \(\ell\). It needs the cited mean-value
theorem for multiplicative functions and the relevant short-interval uniformity.
Track the cancellation of the factors at \(p\le R\) and the replacement of
\(H_{\ell,R}\) by \(H_\ell\); neither can be dropped at the definition stage.

The outputs are Propositions 9.17 and 9.20. Lemma 9.22 bounds the common main
term below. Lemma 9.24 compares the sieve constants. Their combination gives
(9.25), and the checked finite result in Section 4 then applies. The remaining
asymptotic task is to choose the constants in the correct order, show the
\(O(N^{\varepsilon})\) nuisance is smaller than the positive main term, discard
the initial range, and pass from scale \(N^2\) to general \(X\).

\page
\section{Proposed external-input register}
\status{Candidate interfaces; not declarations already installed in the starter}

The user's permission to leave deep theorems unformalized makes a conditional
development appropriate. The exact declarations should be approved mathematically
by comparison with their sources. The list below is an initial register, not a
claim that each result is absent from current Mathlib.

\small
\begin{longtable}{p{.11\linewidth}p{.39\linewidth}p{.41\linewidth}}
\toprule
ID & Published input & Required interface and use\\\midrule
NT01 & PNT in progressions with exceptional term and Landau--Page; paper (2.5)--(2.6), reference [22] &
Uniform growing-modulus error, one datum on an annulus, induced characters, and
uniqueness strong enough across scales. Used in Section 6.\\
NT02 & Siegel--Walfisz and fixed-modulus equidistribution &
Full quantifier order for modulus, residue, exponent and threshold. Used where
polylogarithmic moduli or fixed progressions suffice; not a replacement for NT01.\\
NT03 & Ren [40, Theorem 1] &
Quadratic prime exponential-sum cancellation with exactly the denominator and
arc-width ranges needed in (6.7) and (10.1).\\
NT04 & Bombieri--Vinogradov, Brun--Titchmarsh, Mertens-type estimates; Section 2.6 &
Maxima over residues and endpoints, divisor weights, and constants uniform in
the auxiliary moduli. Used in Section 10 and the sieve.\\
NT05 & Classical quadratic Gauss evaluations and Weil character-sum bounds &
Primitive and induced character conventions; primes and prime powers, including
2; squarefree-polynomial hypothesis for Lemma 8.8.\\
HA01 & Separated-frequency maximal inequality [25, Theorem 1.1] &
Precise separation and cardinality dependence. Prove its adaptation, transference
and uniformity in \(\beta\) in Lemma 6.16.\\
HA02 & Tao [45], Lemma 5.1, Theorems 1.6--1.7 and Proposition 1.5 &
Quantitative major-arc construction, multiplier theorem and sampling comparisons.
Keep \(r,k,B,P,\varepsilon,p\) and vector-valued constants explicit.\\
HA03 & Hardy--Littlewood, Calder\'on--Zygmund, Khintchine and Rademacher--Menshov &
Only use an external declaration where the pinned library does not already
provide the needed statement. Preserve the \(p\downarrow1\) dependence.\\
\bottomrule
\end{longtable}
\normalsize

Do not collapse these into a single assumption saying that all estimates needed
in the paper hold. An import should be independently recognizable as a published
theorem and should not assume the novel conclusion to be checked.

\page
\subsection*{External inputs continued}
\small
\begin{longtable}{p{.11\linewidth}p{.39\linewidth}p{.41\linewidth}}
\toprule
ID & Published input & Required interface and use\\\midrule
ER01 & LaVictoire [27, Theorem 3.1 and Lemma 3.2], with the principles used in the paper's Theorem 8.1 &
CRT families, two coprime sequences, thickening deficit, pair separation and
residue-frequency hypotheses; full persistent universal divergence conclusion.\\
ER02 & Granville--Kurlberg [14, Theorem 1], in the specialization 8.6 &
Uniform correlations for distinct shifts and the hypotheses on local densities;
convergence of normalized cyclic gaps.\\
ER03 & Polynomial-prime pointwise convergence [35, Corollary 2(ii)] &
The bounded-function dense class used in Corollary 7.42, with the correct weights
and systems. Prime powers and Abel summation remain internal tasks.\\
SV01 & Rosser--Iwaniec coefficient-form sieve and fundamental lemma [21] &
Weights, support, signs, coefficient bounds, product hypotheses, sieve functions
and quantitative remainder; prove the two-stage composition in Lemma 11.13.\\
SV02 & Granville--Koukoulopoulos [13] and Shiu [41] &
Mean values of the particular multiplicative weight and uniform estimates in the
retained ranges used by Lemma 11.19.\\\bottomrule
\end{longtable}
\normalsize

\subsection*{Represent imports as hypotheses first}
For a new branch it is often helpful to package the external inputs as a named
structure and prove a theorem with that structure as an explicit argument.
This makes the logical boundary visible without immediately adding global
axioms. Once the interfaces have been audited, a separate file can instantiate
the structure from selected quoted theorems.

In either design, dependency output for a headline theorem must agree with the
published import register. A proof conditional on an estimate proved nowhere in
the project or cited literature should be labeled incomplete, even if Lean
accepts it with that estimate as a hypothesis.

\subsection*{Possible reuse from the Hua project}
The Hua repository contains finite cyclic Fourier infrastructure, quadratic local
arithmetic and analytic assembly. Some may be reusable after comparing
definitions. Its prime-square weights are those of its own \(W\)-trick argument,
not automatically the operators \(A_{N,\Lambda}\) in this paper. Reuse requires
checking signs, normalizations, finite versus infinite groups, and transitive
axiom dependencies. No reuse has been counted as completed work here.

\page
\section{Points to settle during manuscript proofreading}
These are formalization requirements or visible editorial issues, not a verdict
that the main mathematical arguments are wrong.

\begin{enumerate}
\item \textbf{Freeze the version.} The attached paper has 96 pages and Section 7.
Its current numbering is the one used in this document. Match the final LaTeX
labels rather than importing numbering from older discussions.
\item \textbf{Repair the Section 7 references.} The opening of Section 7 and the
last sentence before the proof of Corollary 7.6 refer to Theorem 1.5. In this
snapshot Theorem 1.5 is the divergence theorem, so these references cannot name
the maximal-inequality consequence intended there.
\item \textbf{Preserve the two endpoint functions.} The improving loss is
\(\exp(c\sqrt{L}/L_2)\); the modular function uses
\(\exp(C\sqrt{\log(e+t)\log(e+\log(e+t))})\). These are different scales.
\item \textbf{Separate bounded and unbounded input.} A negligible prime-power
contribution for bounded functions is sufficient for the dense-class convergence
step; it is not itself a uniform error bound for every \(L^1\) function.
\item \textbf{Make exceptional data consistent.} Record the annular data used in
partial summation and prove the regrouping over \(\mathcal M(\chi)\). Do not
replace this by a single exceptional character fixed for all \(N\).
\item \textbf{Expose uniformity in the Goldbach range.} All \(o(1)\), tail errors
and positivity bounds must be uniform in admissible \(n,\ell\) and auxiliary
moduli. Write the order of choosing \(A_0,K,C\) and all thresholds explicitly.
\item \textbf{Retain fixed sieve-prime sets.} The independence from \(n\) after
Lemma 11.2 is required by the mean-square estimate. The coefficient combination
in the pre-sieve must preserve this independence.
\item \textbf{Check strictness and small cases.} Empty sets, \(q=1\), powers of 2,
\(N\ge2\), \(p=2,3\), and equality at density/threshold boundaries should be
covered by exact definitions, not silently discarded.
\end{enumerate}

The present work has not audited every cited theorem or every proof in the PDF.
In particular, successful checks of the starter do not validate any unfinished
node elsewhere in the map.

\page
\section{What is needed to continue}
The PDFs were enough to start. The next useful input is the final LaTeX project:
all source files, bibliography, custom macros and any code used for numerical
estimates, together with the PDF generated from those files. This makes the
statements and labels unambiguous and lets the map track proofreading changes.

GitHub is useful for versioning and automated builds, but an account is not
required to write or check Lean locally. This archive is a normal Lean project;
it can be uploaded later. The toolchain and lockfile are already selected.
There is no need for the author to learn Lean before mathematical work can begin.

The next mathematical milestone should be one coherent branch. A practical
sequence is:
\begin{enumerate}
\item Encode the final target statements and audit their quantifiers against the
LaTeX. Build a stable statement map.
\item Complete a substantial finite arithmetic module, such as the quadratic
Cohen identity with all prime-power cases, or complete the scalar optimization
in Theorem 4.1.
\item Specify and audit the deep-input interfaces. Prove one full branch from
those interfaces, without assuming its novel intermediate estimates.
\item Add the remaining branches and report the exact dependencies of each
headline theorem separately.
\end{enumerate}

For the Goldbach branch, the existing checked finite assembly supplies an endpoint
to which the circle-method and sieve components can connect. For the harmonic-
analysis branch, the shared high--low decomposition is the principal reusable
milestone. The divergence branch can be developed separately once its external
criterion and Poisson-spacings interfaces are audited.

\textbf{Completion criterion.} A full conditional formalization requires a fresh
build of every target theorem, a transitive axiom audit with no unfinished proof
axiom, and a human comparison of the final statements and imported inputs with
the manuscript and sources. A long blueprint, a clean source scan, or a passing
build of the present starter alone does not meet that criterion.

\textbf{Current deliverable.} This is an initial, reproducible starting point with
small checked components and an explicit account of the much larger unfinished
work. It should not yet be described in an arXiv or journal submission as a
formalization of the paper.

\page
\appendix
\section{Inventory of numbered statements}
The following inventory was extracted from headings in the attached PDF. Planned
names are suggestions only. A status of ``target'' means no proof of that numbered
statement is present in the starter. Checked supporting subclaims are described
separately in Sections 4--5; they do not change the status of the enclosing paper
theorem or lemma.
\small
\begin{longtable}{p{.23\linewidth}r p{.53\linewidth}p{.1\linewidth}}
\toprule Statement & Page & Planned name (after \texttt{PrimeSquares.}) & Status\\\midrule
\endhead
Theorem 1.1 & 3 & \lean{Main.theorem_1_1} & target\\
Theorem 1.3 & 3 & \lean{Main.theorem_1_3} & target\\
Theorem 1.5 & 4 & \lean{Main.theorem_1_5} & target\\
Theorem 1.6 & 5 & \lean{Main.theorem_1_6} & target\\
Definition 2.10 & 12 & \lean{Foundations.definition_2_10} & target\\
Definition 2.11 & 13 & \lean{Foundations.definition_2_11} & target\\
Lemma 3.1 & 15 & \lean{Arithmetic.lemma_3_1} & target\\
Definition 3.7 & 16 & \lean{Arithmetic.definition_3_7} & target\\
Definition 3.9 & 17 & \lean{Arithmetic.definition_3_9} & target\\
Lemma 3.11 & 17 & \lean{Arithmetic.lemma_3_11} & target\\
Lemma 3.12 & 17 & \lean{Arithmetic.lemma_3_12} & target\\
Lemma 3.14 & 17 & \lean{Arithmetic.lemma_3_14} & target\\
Corollary 3.16 & 18 & \lean{Arithmetic.corollary_3_16} & target\\
Lemma 3.20 & 19 & \lean{Arithmetic.lemma_3_20} & target\\
Lemma 3.21 & 19 & \lean{Arithmetic.lemma_3_21} & target\\
Lemma 3.24 & 20 & \lean{Arithmetic.lemma_3_24} & target\\
Lemma 3.27 & 21 & \lean{Arithmetic.lemma_3_27} & target\\
Theorem 4.1 & 24 & \lean{Improving.theorem_4_1} & target\\
Corollary 4.15 & 25 & \lean{Improving.corollary_4_15} & target\\
Proposition 4.19 & 26 & \lean{Improving.proposition_4_19} & target\\
Lemma 5.1 & 28 & \lean{Sparse.lemma_5_1} & target\\
Lemma 5.10 & 30 & \lean{Sparse.lemma_5_10} & target\\
Lemma 5.17 & 31 & \lean{Sparse.lemma_5_17} & target\\
Lemma 6.5 & 34 & \lean{HighLow.lemma_6_5} & target\\
Lemma 6.9 & 35 & \lean{HighLow.lemma_6_9} & target\\
Proposition 6.14 & 37 & \lean{HighLow.proposition_6_14} & target\\
Lemma 6.16 & 39 & \lean{HighLow.lemma_6_16} & target\\
Corollary 6.19 & 40 & \lean{HighLow.corollary_6_19} & target\\
Proposition 6.21 & 41 & \lean{HighLow.proposition_6_21} & target\\
Theorem 7.2 & 44 & \lean{Maximal.theorem_7_2} & target\\
Corollary 7.6 & 44 & \lean{Maximal.corollary_7_6} & target\\
Proposition 7.8 & 45 & \lean{Maximal.proposition_7_8} & target\\
Lemma 7.15 & 47 & \lean{Maximal.lemma_7_15} & target\\
Proposition 7.21 & 48 & \lean{Maximal.proposition_7_21} & target\\
Proposition 7.28 & 51 & \lean{Maximal.proposition_7_28} & target\\
Corollary 7.42 & 56 & \lean{Maximal.corollary_7_42} & target\\
Theorem 8.1 & 57 & \lean{Divergence.theorem_8_1} & target\\
Theorem 8.6 & 58 & \lean{Divergence.theorem_8_6} & target\\
Lemma 8.8 & 58 & \lean{Divergence.lemma_8_8} & target\\
Lemma 9.2 & 61 & \lean{Goldbach.lemma_9_2} & target\\
Lemma 9.3 & 61 & \lean{Goldbach.lemma_9_3} & target\\
Lemma 9.7 & 62 & \lean{Goldbach.lemma_9_7} & target\\
Corollary 9.9 & 63 & \lean{Goldbach.corollary_9_9} & target\\
Lemma 9.10 & 63 & \lean{Goldbach.lemma_9_10} & target\\
Proposition 9.17 & 64 & \lean{Goldbach.proposition_9_17} & target\\
Proposition 9.20 & 64 & \lean{Goldbach.proposition_9_20} & target\\
Lemma 9.22 & 65 & \lean{Goldbach.lemma_9_22} & target\\
Lemma 9.24 & 65 & \lean{Goldbach.lemma_9_24} & target\\
Lemma 10.2 & 66 & \lean{GoldbachCircle.lemma_10_2} & target\\
Lemma 10.12 & 68 & \lean{GoldbachCircle.lemma_10_12} & target\\
Lemma 10.15 & 68 & \lean{GoldbachCircle.lemma_10_15} & target\\
Lemma 10.20 & 72 & \lean{GoldbachCircle.lemma_10_20} & target\\
Lemma 10.24 & 73 & \lean{GoldbachCircle.lemma_10_24} & target\\
Lemma 11.2 & 76 & \lean{GoldbachSieve.lemma_11_2} & target\\
Lemma 11.6 & 76 & \lean{GoldbachSieve.lemma_11_6} & target\\
Lemma 11.8 & 77 & \lean{GoldbachSieve.lemma_11_8} & target\\
Lemma 11.13 & 78 & \lean{GoldbachSieve.lemma_11_13} & target\\
Lemma 11.19 & 82 & \lean{GoldbachSieve.lemma_11_19} & target\\
\bottomrule
\end{longtable}

\normalsize
\end{document}
